\begin{abstract}
Style representations aim to embed texts with similar writing styles closely and texts with different styles far apart, regardless of content. However, the contrastive triplets often used for training these representations may vary in both style and content, leading to potential content leakage in the representations. We introduce \textsc{StyleDistance}, a novel approach to training stronger content-independent style embeddings. We use a large language model to create a synthetic dataset of near-exact paraphrases with controlled style variations, and produce positive and negative examples across 40 distinct style features for precise contrastive learning. We assess the quality of our synthetic data and embeddings through human and automatic evaluations.  \textsc{StyleDistance} enhances the content-independence of style embeddings, which generalize to real-world benchmarks and outperform leading style representations in downstream applications. Our model can be found at \url{https://huggingface.co/StyleDistance/styledistance}.
% \kmnote{Abstract is a bit too long}
%Style representations aim to embed texts with similar writing styles closely in an embedding space and texts with different styles far apart, regardless of content.  \kmnote{A couple of problems in the next 3 sentences: too much on prior work. Not appropriate for abstract. Also, "remains unclear" is weak. One fix: Cut next 3 sentences and replace with "However, the contrastive triplets used for training these representations may vary by both style and content, leading to potential content leakage." you really should have the focus on your contribution so get to "in this paper" quicker.} However, the extent of content-independence in current style embeddings remains unclear.  Existing models often train using contrastive triplets constructed from social media datasets, with a proxy objective that uses texts by the same author as positive examples and texts by different authors as negative examples. These triplets, however, are not perfect contrastive examples in terms of style (i.e. both style and content might vary), leading to potential content leakage. In this paper, we introduce \textsc{StyleDistance}, a novel approach to training stronger content-independent style embeddings. We use a large language model to create a synthetic dataset of near-exact paraphrases with controlled style variations, and produce positive and negative examples across 40 distinct style features for precise contrastive learning. We assess the quality of our synthetic data and embeddings through human and automated evaluations. We evaluate the quality of our embeddings on the STEL and STEL-or-Content tasks which were designed to measure the quality of style representations and their content-independence. We find training on these synthetic parallel examples significantly enhances the content-independence of style embeddings, which generalize to real-world benchmarks and outperform leading style representations in downstream applications.
\end{abstract}
\mainfig



